\documentclass[fleqn]{beamer}

\mode<presentation> {
\usetheme{Madrid}
\usecolortheme{dolphin}
}

\usepackage[utf8]{inputenc}
\usepackage[brazilian]{babel}
\usepackage{amsmath}
\usepackage{amssymb}
\usepackage{booktabs}
\usepackage{listings}
\usepackage{xcolor}

\lstset{
	basicstyle=\ttfamily\scriptsize,
	breaklines=true,
	keywordstyle=\color{blue!60!black},
	commentstyle=\color{gray},
	stringstyle=\color{green!45!black},
	showstringspaces=false,
	columns=fullflexible,
	frame=single,
	framesep=4pt,
	aboveskip=6pt,
	belowskip=2pt,
}

\title[Aula 3 -- AMPL + Python]{Transbordo e Localização de Facilidades}
\subtitle{Rodando o AMPL de dentro do Python}
\author[Santi]{Prof. \'Everton Santi}
\institute[UFRN/ECT]{ECT2524 -- Introdução à Otimização \\ Escola de Ciências e Tecnologia -- UFRN}
\date{Aula 3}

\begin{document}

\begin{frame}
	\titlepage
\end{frame}

\begin{frame}
	\frametitle{Sumário}
	\tableofcontents
\end{frame}

%------------------------------------------------------------
\section{Recapitulando}
%------------------------------------------------------------

\begin{frame}
	\frametitle{Até aqui}
	\begin{itemize}
		\item Aulas 1 e 2: modelar $\rightarrow$ escrever em AMPL
			(\texttt{.mod}, \texttt{.dat}, \texttt{.run}) $\rightarrow$
			resolver com HiGHS pela linha de comando;
		\item Isso funciona bem para \textbf{um} conjunto de dados escrito
			à mão no \texttt{.dat};
		\item Mas e quando temos \textbf{dezenas de instâncias}, em um
			formato de arquivo que não é \texttt{.dat}, e queremos rodar
			todas e comparar?
	\end{itemize}
\end{frame}

\begin{frame}
	\frametitle{Hoje: AMPL chamado de dentro do Python}
	\begin{itemize}
		\item Mesmo AMPL, mesmo solver -- só que dirigido por um
			\textbf{script Python} (biblioteca \texttt{amplpy});
		\item Ganhos:
		\begin{itemize}
			\item ler dados de \textbf{qualquer formato} (\texttt{.txt},
				\texttt{.csv}, planilha\dots) com código Python;
			\item montar os parâmetros do modelo \textbf{por programa},
				não copiando número a número;
			\item rodar \textbf{um laço} sobre várias instâncias e
				coletar resultados numa tabela;
		\end{itemize}
		\item A formulação e o modelo vamos construir \textbf{ao vivo},
			em Python, durante a aula;
		\item Estes slides trazem só a \textbf{descrição do problema} e a
			\textbf{colinha de comandos} do \texttt{amplpy}.
	\end{itemize}
\end{frame}

%------------------------------------------------------------
\section{O Problema de Transbordo}
%------------------------------------------------------------

\begin{frame}
	\frametitle{O contexto}
	\begin{itemize}
		\item Uma empresa do ramo de \textbf{bebidas} precisa organizar a
			distribuição para o feriado de \textbf{carnaval};
		\item A rede tem \textbf{três camadas}:
		\begin{itemize}
			\item \textbf{fábricas} produzem o produto;
			\item \textbf{centros de distribuição} recebem das fábricas e
				repassam;
			\item \textbf{clientes} têm uma demanda a ser atendida;
		\end{itemize}
		\item O produto só chega ao cliente \textbf{passando por um
			centro} -- daí o nome \emph{transbordo} (troca de veículo /
			ponto intermediário).
	\end{itemize}
\end{frame}

\begin{frame}
	\frametitle{Custos envolvidos}
	\begin{itemize}
		\item \textbf{Transporte} (custo variável, por unidade):
		\begin{itemize}
			\item $c_{ij}$: produzir na fábrica $i$ e levar ao centro $j$;
			\item $t_{jk}$: armazenar no centro $j$ e entregar ao cliente $k$;
		\end{itemize}
		\item \textbf{Custo fixo} (pago uma vez, só se a instalação for usada):
		\begin{itemize}
			\item $f_i$: abrir/operar a fábrica $i$;
			\item $h_j$: abrir/operar o centro $j$;
		\end{itemize}
		\item \textbf{Capacidade} (limite de quanto passa pela instalação):
		\begin{itemize}
			\item $g_i$: capacidade da fábrica $i$;
			\item $q_j$: capacidade do centro $j$;
		\end{itemize}
		\item \textbf{Demanda}: $d_k$ unidades exigidas pelo cliente $k$.
	\end{itemize}
\end{frame}

\begin{frame}
	\frametitle{O que precisamos decidir}
	\begin{itemize}
		\item Quais \textbf{fábricas} abrir e quais \textbf{centros} abrir;
		\item Quanto produto enviar em \textbf{cada trecho}
			fábrica~$\rightarrow$~centro e centro~$\rightarrow$~cliente;
		\item Regras que a solução tem que respeitar:
		\begin{itemize}
			\item nenhuma instalação passa mais do que sua capacidade;
			\item a demanda de cada cliente é atendida \textbf{exatamente};
			\item \textbf{conservação de fluxo} no centro: tudo que entra
				sai (o centro não acumula estoque);
			\item só há fluxo saindo de uma instalação se ela estiver
				\textbf{aberta}.
		\end{itemize}
		\item \textbf{Objetivo}: minimizar o custo total (fixos + variáveis).
	\end{itemize}
\end{frame}

\begin{frame}
	\frametitle{Dados da instância do exemplo}
	\begin{columns}[t]
		\column{0.5\textwidth}
		2 fábricas, 3 centros, 3 clientes.
		\medskip

		Demanda dos clientes:
		\begin{center}
		\begin{tabular}{cccc}
			\toprule
			$k$ & 1 & 2 & 3 \\
			$d_k$ & 80 & 90 & 40 \\
			\bottomrule
		\end{tabular}
		\end{center}

		Fábricas:
		\begin{center}
		\begin{tabular}{ccc}
			\toprule
			 & $i{=}1$ & $i{=}2$ \\
			\midrule
			$g_i$ & 120 & 100 \\
			$f_i$ & 300 & 250 \\
			\bottomrule
		\end{tabular}
		\end{center}

		Centros:
		\begin{center}
		\begin{tabular}{cccc}
			\toprule
			 & $j{=}1$ & $j{=}2$ & $j{=}3$ \\
			\midrule
			$q_j$ & 70 & 90 & 66 \\
			$h_j$ & 60 & 80 & 100 \\
			\bottomrule
		\end{tabular}
		\end{center}

		\column{0.5\textwidth}
		Custo fábrica $\rightarrow$ centro $c_{ij}$:
		\begin{center}
		\begin{tabular}{cccc}
			\toprule
			$c_{ij}$ & $j{=}1$ & $j{=}2$ & $j{=}3$ \\
			\midrule
			$i{=}1$ & 10 & 8 & 20 \\
			$i{=}2$ & 20 & 6 & 11 \\
			\bottomrule
		\end{tabular}
		\end{center}
		\medskip

		Custo centro $\rightarrow$ cliente $t_{jk}$:
		\begin{center}
		\begin{tabular}{cccc}
			\toprule
			$t_{jk}$ & $k{=}1$ & $k{=}2$ & $k{=}3$ \\
			\midrule
			$j{=}1$ & 7 & 5 & 10 \\
			$j{=}2$ & 5 & 4 & 3 \\
			$j{=}3$ & 10 & 8 & 6 \\
			\bottomrule
		\end{tabular}
		\end{center}
		\medskip

		\footnotesize Demanda total $= 210$; capacidade das fábricas
		$= 220$; dos centros $= 226$.
	\end{columns}
\end{frame}

%------------------------------------------------------------
\section{O arquivo de dados}
%------------------------------------------------------------

\begin{frame}[fragile]
	\frametitle{Formato do \texttt{transbordo.txt}}
	\begin{lstlisting}
2 3 3
80 90 40
120 100
300 250
70 90 66
60 80 100
10 8 20
20 6 11
7 5 10
5 4 3
10 8 6
	\end{lstlisting}
	\begin{itemize}
		\item linha 1: \texttt{m n p} -- nº de fábricas, centros, clientes;
		\item linha 2: demandas $d_k$ ($p$ valores);
		\item linhas 3--4: $g_i$ e $f_i$ das fábricas ($m$ valores cada);
		\item linhas 5--6: $q_j$ e $h_j$ dos centros ($n$ valores cada);
		\item próximas $m$ linhas: matriz $c_{ij}$ ($n$ colunas);
		\item próximas $n$ linhas: matriz $t_{jk}$ ($p$ colunas).
	\end{itemize}
\end{frame}

%------------------------------------------------------------
\section{amplpy: a colinha}
%------------------------------------------------------------

\begin{frame}[fragile]
	\frametitle{Preparar o ambiente}
	\begin{lstlisting}[language=Python]
# instalar uma vez no ambiente
pip install amplpy

# ativar a licenca e os modulos de solver (uma vez)
python -m amplpy.modules install highs
python -m amplpy.modules activate <sua-licenca-uuid>
	\end{lstlisting}
	\begin{itemize}
		\item \texttt{amplpy} = ponte Python $\leftrightarrow$ AMPL;
		\item o módulo \texttt{highs} traz o solver HiGHS (mesmo das
			aulas anteriores);
		\item a licença acadêmica é gratuita em \url{https://ampl.com}.
	\end{itemize}
\end{frame}

\begin{frame}[fragile]
	\frametitle{Criar o modelo e passar os dados}
	\begin{lstlisting}[language=Python]
from amplpy import AMPL

ampl = AMPL()

# opcao A: escrever o modelo inline
ampl.eval(r"""
    param m > 0;
    param d{1..m} >= 0;
    var x{1..m} >= 0;
    minimize custo: sum{i in 1..m} c[i]*x[i];
    ...
""")
# opcao B: ler de um arquivo .mod
ampl.read("transbordo.mod")

# parametro escalar
ampl.param["m"] = dados["m"]

# vetor  -> dict {indice: valor}
ampl.param["d"] = {k+1: v for k, v in enumerate(dados["d"])}

# matriz -> dict {(i, j): valor}
ampl.param["c"] = {(i+1, j+1): dados["c"][i][j]
                   for i in range(m) for j in range(n)}
	\end{lstlisting}
\end{frame}

\begin{frame}[fragile]
	\frametitle{Resolver e ler a solução}
	\begin{lstlisting}[language=Python]
ampl.set_option("solver", "highs")
ampl.set_option("highs_options", "timelim=60 outlev=1")

ampl.solve()

# status: 'solved', 'limit', 'infeasible', ...
print(ampl.solve_result)

# valor da funcao objetivo
custo = ampl.get_objective("custo").value()

# valores das variaveis
x = ampl.get_variable("x").get_values().to_pandas()
y = ampl.get_variable("y").get_values().to_list()

# limitante inferior / gap (quando o solver expoe)
lb = ampl.get_value("custo.bestbound")
gap = ampl.get_value("custo.relmipgap")
	\end{lstlisting}
	\begin{itemize}
		\item para rodar várias instâncias: pôr tudo numa função e
			chamar num laço sobre a lista de arquivos.
	\end{itemize}
\end{frame}

%------------------------------------------------------------
\section{Roteiro do exemplo}
%------------------------------------------------------------

\begin{frame}
	\frametitle{O que vamos fazer ao vivo}
	\begin{enumerate}
		\item Escrever \texttt{ler\_dados(arquivo)}: abre o
			\texttt{transbordo.txt} e devolve um dicionário
			(\texttt{m, n, p, d, g, f, q, h, c, t});
		\item Escrever o \textbf{modelo AMPL} do transbordo (variáveis de
			abertura, variáveis de fluxo, objetivo, restrições);
		\item Escrever o \textbf{script Python}: lê os dados, injeta nos
			parâmetros, resolve;
		\item \textbf{Ler a solução}: quais fábricas e centros abriram,
			quanto passa em cada trecho, custo total;
		\item Discutir: como isso escala para \textbf{muitas instâncias}
			e para \textbf{experimentos} (tempo limite, gap, comparação).
	\end{enumerate}
\end{frame}

%------------------------------------------------------------
\section{Tarefa}
%------------------------------------------------------------

\begin{frame}
	\frametitle{Tarefa 3 -- Localização de Facilidades}
	\begin{itemize}
		\item Estudo de caso: escolher \textbf{locais de prova} para um
			exame nacional -- há candidatos, locais candidatos com
			\textbf{capacidade}, um limite de quantos locais podem ser
			usados, e a \textbf{distância} de cada candidato a cada local;
		\item Sigam \textbf{os mesmos passos do exemplo de hoje}:
		\begin{enumerate}
			\item \textbf{modelar} -- vocês identificam parâmetros,
				variáveis, objetivo e restrições (o enunciado dá
				\textbf{duas opções de objetivo}: soma das distâncias ou
				a maior distância -- escolham e justifiquem);
			\item escrever o \textbf{modelo em AMPL}
				(\texttt{facilidades.mod});
			\item escrever o \textbf{código Python}: \texttt{ler\_dados}
				+ script \texttt{amplpy} que resolve as \textbf{3
				instâncias} e monta a tabela de resultados (objetivo, LB,
				GAP, tempo, status).
		\end{enumerate}
		\item Enunciado completo e arquivos \texttt{.txt} na página da
			aula, no site do curso.
	\end{itemize}
\end{frame}

\begin{frame}
	\frametitle{Referências}
	\begin{itemize}
		\item AMPL: \url{https://ampl.com};
		\item Documentação do \texttt{amplpy}: \url{https://amplpy.ampl.com}.
	\end{itemize}
\end{frame}

\end{document}
